% Generated draft: requires mathematical and source review.
\documentclass[11pt]{article}
\usepackage[margin=25mm]{geometry}
\usepackage{fontspec}
\setmainfont{Latin Modern Roman}
\usepackage{amsmath,amssymb,mathtools,unicode-math}
\setmathfont{Latin Modern Math}
\usepackage{enumitem,needspace,xurl,hyperref}
\setlength{\parindent}{0pt}
\setlength{\parskip}{5pt}
\setlength{\emergencystretch}{3em}
\newcommand{\N}{\mathbb{N}}
\newcommand{\Z}{\mathbb{Z}}
\newcommand{\Q}{\mathbb{Q}}
\newcommand{\R}{\mathbb{R}}
\newcommand{\C}{\mathbb{C}}
\newcommand{\F}{\mathbb{F}}
\newcommand{\A}{\mathbb{A}}
\newcommand{\PP}{\mathbb{P}}
\newcommand{\E}{\mathbb{E}}
\newcommand{\eps}{\varepsilon}
\newcommand{\dd}{\,\mathrm{d}}
\begin{document}
\section*{UnsolvedMath ID: 700880}
\textbf{Source problem code:} catalogue entry
\subsection{Definitions and mathematical statement}
Let \(\Omega\) be an infinite set and let \(G\leq\operatorname{Sym}(\Omega)\) act transitively on \(\Omega\). For \(\alpha\in\Omega\), define the point stabilizer by
\[
G_\alpha=\{g\in G:g\alpha=\alpha\}.
\]
Following the definition in the supplied record, call this permutation group a Frobenius group if
\[
G_\alpha\cap G_\beta=\{1\}\qquad\text{for all distinct }\alpha,\beta\in\Omega.
\]
Thus the stabilizer of two points means their pointwise stabilizer. The source definition does not separately require a nontrivial point stabilizer.

The action is primitive if its only \(G\)-invariant equivalence relations on \(\Omega\) are equality and the universal relation. For a subgroup \(H\leq G\) and \(\beta\in\Omega\), write \(H\beta=\{h\beta:h\in H\}\) for the corresponding orbit.

\textbf{Existence question (P. M. Neumann, Kourovka Notebook Problem 12.62).} Does there exist such a Frobenius group of infinite degree whose action is primitive and for which, for some \(\alpha\in\Omega\),
\[
G_\alpha=\langle c\rangle\quad\text{for some }c\in G_\alpha,
\qquad
\bigl|\{G_\alpha\beta:\beta\in\Omega\}\bigr|<\infty?
\]
By transitivity, requiring these properties at one point is equivalent to requiring them at every point. The orbit count is on all of \(\Omega\), including the singleton orbit \(\{\alpha\}\); cyclicity carries no additional assumption that the cyclic group is finite or infinite.

\textbf{Supplied conjecture.} No such permutation group exists.

\textbf{Supplied qualification.} The record states that any group satisfying the proposed conditions cannot be \(2\)-transitive, meaning transitive on the ordered pairs of distinct points of \(\Omega\), and gives the references listed below. This qualification is retained as a source assertion and has not been independently verified here. The registry labels the problem open; this is a supplied status, not an independent literature assessment.
\subsection{Short English statement}
P. M. Neumann asks whether an infinite primitive transitive permutation group can have trivial stabilizers of pairs of distinct points and a cyclic point stabilizer with only finitely many orbits. The supplied conjecture says no. The source also states that such a group cannot be \(2\)-transitive.
\subsection{Sources}
\begin{itemize}
\item Local UnsolvedMath registry, permanent ID \(700880\), exact source code \texttt{catalogue entry}, title \emph{Kourovka Notebook Problem 12.62}. The supplied record attributes the problem to P. M. Neumann, with proposed year 1992, and labels its status open. Its recorded update date is \texttt{2026-10-16}; this is registry metadata, not evidence of independent source verification.
\item \emph{The Kourovka Notebook (Unsolved Problems in Group Theory)}, Problem 12.62, 12th issue (1992); the supplied background also identifies the 21st issue edition (2026), arXiv:1401.0300v48. The collection metadata names E. I. Khukhro and V. D. Mazurov as editors. These source identifications are supplied attribution and have not been independently checked here.
\item Gy. K\'arolyi, S. J. Kov\'acs, P. P. P\'alfy, \emph{Aequationes Math.}, \textbf{39} (1990), 161--166. Cited in the supplied statement in support of the qualification concerning \(2\)-transitivity.
\item V. D. Mazurov, \emph{Siberian Math. J.}, \textbf{31} (1990), 615--617. Cited in the supplied statement in support of the qualification concerning \(2\)-transitivity.
\item P. M. Neumann, P. J. Rowley, in: \emph{Geometry and cohomology in group theory}, London Math. Soc. Lect. Note Ser., \textbf{252}, Cambridge Univ. Press, 1998, 291--295. Cited in the supplied statement in support of the qualification concerning \(2\)-transitivity.
\end{itemize}
\end{document}
